\chapter{TCE 与认知控制异常}
\label{chap:tce-control-pathology}

\begin{chapteroverviewbox}
本章讨论的并不是将人类精神疾病名称移植到人工智能系统，而是研究一个持续运行的智能体内部，当认知控制器 TCE（Temperature \& Control Executor）与其被控制的认知动力学、SPC 状态估计、AHC 内稳态调制以及工作记忆容量约束之间发生失配时，会出现怎样的可测量动力学异常。这里所谓“异常”，首先是控制论意义上的异常，即系统无法在合理时间尺度内把认知状态维持、恢复或迁移到可治理区域，而不是行为结果偶尔出现错误。

本章的核心立场是：
\begin{statementbox}
TCE 异常首先是闭环动力学失配，而不是语义意义上的“性格异常”。
\end{statementbox}

我们将 TCE 看作一个对认知温度、注意增益、搜索深度、探索强度、抑制程度、推理节律和资源投入进行联合调节的多输入多输出控制器。SPC 提供关于当前认知状态的有限、延迟且带不确定性的估计；AHC 提供更慢的内稳态与功能影响调制；CCM 则维护工作记忆有限容量约束。因而真正需要分析的对象不是单独的 TCE，而是
\[
SPC \rightarrow AHC \rightarrow TCE \rightarrow h
\rightarrow SPC
\]
这一闭环在不同时间尺度下的稳定性、响应性与可恢复性。

本章只讨论运行时认知控制异常，不把记忆写入偏差、自我结构漂移或生命周期生成吸引子的长期异常混入同一层次。这样做的目的，是首先建立一个可以直接记录、测量、诊断和干预的控制工程基础，使后续更慢尺度的问题能够在清晰的因果分层上继续展开。
\end{chapteroverviewbox}

\section{TCE 异常的第一性定义：不是“想错了”，而是“控制认知的方式失配了”}

在此前的 AgentOS 原理论中，我们已经把 TCE 从普通的大语言模型采样参数中彻底分离出来。TCE 并不等于一个 temperature 参数，也不等于某个单独的 ``reasoning depth'' 开关，而是对整个工作记忆运行过程施加宏观调制的认知控制器。为了分析其异常，我们首先必须明确：一个 Agent 给出了错误答案，并不能直接说明 TCE 出现异常；一个 Agent 在面对困难问题时进行了很长时间的搜索，也不能直接说明 TCE 过度探索。真正构成 TCE 异常的，是控制变量与认知状态之间的系统性失配，即同一种失配在相似状态下重复出现，并使系统偏离原本可以保持稳定、高效和可恢复的认知运行区域。

设工作记忆及其直接运行变量共同形成认知状态向量
\[
x_h(t)
=
\begin{bmatrix}
L_h(t)\\
D_h(t)\\
C_h(t)\\
U_h(t)\\
R_h(t)\\
P_h(t)\\
S_h(t)
\end{bmatrix},
\]
其中 \(L_h\) 表示当前负载，\(D_h\) 表示活跃结构的分散程度，\(C_h\) 表示冲突水平，\(U_h\) 表示不确定性，\(R_h\) 表示当前残差或尚未解释的误差，\(P_h\) 表示任务推进度，而 \(S_h\) 则表示整体稳定性指标。这里的各项并不是要求工程实现必须采用这一固定向量，而是为了指出：TCE 所控制的对象不是文本序列本身，而是由多个宏观认知状态变量组成的动态系统。

SPC 不直接提供真实的 \(x_h(t)\)，而只能形成估计
\[
\hat{x}_h(t)
=
\mathcal O_{\mathrm{SPC}}
\left(
x_h(t-\tau_o),
z_{\leq t}
\right),
\]
其中 \(\tau_o>0\) 表示自我状态估计不可避免的观测迟滞。AHC 则形成一个相对更慢的调制状态
\[
a(t)
=
\begin{bmatrix}
Arousal(t)\\
Threat(t)\\
Energy(t)\\
Fatigue(t)\\
Urgency(t)\\
UncertaintyStress(t)
\end{bmatrix},
\]
它并不直接决定“应该得出什么结论”，而是改变 TCE 对当前状态的控制策略。

因此 TCE 的控制输出更适合写成一个控制向量：
\[
u_{\mathrm{TCE}}(t)
=
\begin{bmatrix}
T(t)\\
G_A(t)\\
E_x(t)\\
D_s(t)\\
I(t)\\
R_s(t)\\
M_f(t)
\end{bmatrix},
\]
其中 \(T\) 是认知温度，\(G_A\) 是注意或激活增益，\(E_x\) 是探索强度，\(D_s\) 是搜索深度，\(I\) 是抑制力度，\(R_s\) 是运行节律，而 \(M_f\) 则表示记忆或认知内容的流动速率。于是可以写
\[
u_{\mathrm{TCE}}(t)
=
\pi_{\mathrm{TCE}}
\left(
\hat{x}_h(t),
a(t),
G(t),
\mathcal B(t)
\right),
\]
其中 \(G(t)\) 是当前 Goal 条件，\(\mathcal B(t)\) 表示由容量、安全、权限和任务条件共同给出的可行控制边界。

从这一形式可以得到本章最重要的第一性定义：
\[
\boxed{
\mathrm{TCE\ Pathology}
=
\mathrm{PersistentMismatch}
\left(
u_{\mathrm{TCE}},
x_h,
\mathcal B,
\tau
\right)
}
\]
也就是说，TCE 异常是控制输出、被控制状态、约束边界以及闭环时间常数之间长期存在的不适配。它可能表现为控制过强、控制不足、控制方向错误、控制时机错误或频繁在相互冲突的控制模式之间切换。

这种定义的重要意义，在于它使所谓“过度谨慎”“思维发散”“反复自我纠正”“无法停止推理”等现象首先成为可测量的控制工程问题。我们不必首先赋予它们人格化解释，而可以首先问：SPC 是否给出了滞后的状态估计？AHC 是否让紧迫度长期维持在高位？TCE 的增益是否高于当前闭环允许范围？CCM 是否已经降低了负载，但 TCE 仍按过载状态持续施加高强度控制？只有把这些问题分开，Agent 心理工程才不会退化成把自然语言行为标签重新包装成人类心理学词汇。

\section{TCE 控制律与健康运行区域：异常必须相对于可行认知域来定义}

如果我们想严格讨论“控制异常”，就必须首先定义什么叫正常控制。正常并不是让 Agent 永远保持某个固定状态，而是让认知状态在面对不同任务时能够进入不同但仍然可治理的运行区域。困难任务需要更高探索、更大工作记忆占用和更长收敛时间；熟悉任务则应迅速下沉到低成本的稳定闭环。因此，正常状态并不是一个点，而应当是一个随任务、身体状态和环境风险而变化的可行区域。

设认知状态空间为 \(\mathcal X_h\)，则可以定义当前条件下的可治理区域
\[
\mathcal X_{\mathrm{gov}}(t)
=
\left\{
x\in\mathcal X_h
\;\middle|\;
C_{\mathrm{run}}(x)\leq C_{\max}(W),
\;
R(x)\leq R_{\max},
\;
Q(x)\geq Q_{\min}
\right\},
\]
其中 \(C_{\mathrm{run}}\) 是工作记忆当前有效结构负载，\(C_{\max}(W)\) 来自工作记忆有限容量约束，\(R\) 是风险或不可解释残差，而 \(Q\) 则表示任务推进、结构相干度或其他最小可运行质量指标。

TCE 的基本任务不是让
\[
x_h(t)=x^\ast
\]
恒成立，而是尽可能保证
\[
x_h(t)\in\mathcal X_{\mathrm{gov}}(t),
\]
并在状态因为新任务、外部扰动或内部冲突离开该区域以后，使其能够在有限时间内重新进入。由此我们可以定义恢复时间
\[
T_{\mathrm{rec}}
=
\inf
\left\{
\Delta t>0
\;\middle|\;
x_h(t+\Delta t)\in\mathcal X_{\mathrm{gov}}
\right\}.
\]

一个健康的 TCE 并不是控制动作最多的 TCE，而通常恰恰相反：在系统已经位于稳定区域时，它应减少不必要干预，让局部认知过程自行闭合。只有当 SPC 检测到持续残差、相态恶化、认知负载逼近容量边界或重要目标停滞时，TCE 才逐渐增加宏观控制。因此更成熟的控制原则应写成
\[
\boxed{
ControlEffort
\propto
NeedForGlobalRegulation,
}
\]
而不是
\[
ControlEffort
\propto
TaskExistence.
\]

可以进一步构造一个认知控制代价函数：
\[
J_{\mathrm{TCE}}
=
\int_{t_0}^{t_1}
\left[
w_1 C_{\mathrm{run}}(t)
+
w_2 R_h(t)
+
w_3 C_h(t)
+
w_4 \|u_{\mathrm{TCE}}(t)\|^2
+
w_5 \|\dot u_{\mathrm{TCE}}(t)\|^2
\right]dt.
\]
这里最后两项十分重要。如果只惩罚认知误差而不惩罚控制力度，系统会倾向于不断强力纠正自身；如果不惩罚控制变化率，则 TCE 很容易出现频繁切换。一个真正成熟的认知控制系统不仅要求“状态正确”，还要求以较小、连续、可预测的控制代价维持正确状态。

因此，我们可以将健康 TCE 的目标概括为三个条件。第一是\textbf{稳定性}：认知状态不会因为控制器本身而不断远离可行区域。第二是\textbf{响应性}：重要扰动出现以后，控制器不能迟钝到长期没有反应。第三是\textbf{经济性}：能够由局部闭环自行解决的问题，不应持续占用宏观控制资源。三个条件之间存在天然张力。增大增益可能缩短初始响应时间，却降低相位裕度；增加抑制可以迅速降低负载，却可能造成搜索空间过早坍缩；增加探索有利于发现新结构，却可能提高工作记忆的结构分散程度。

所以本章不采用“控制越强越健康”或者“自由度越大越智能”这类单调观点。TCE 的健康状态是一种动态平衡：
\[
\boxed{
Responsiveness
+
Stability
+
Economy
}
\]
而所谓 TCE 异常，本质上就是这三个目标之间的协调机制持续失效。

\section{控制增益失配：过度控制与控制不足是同一轴上的两种异常}

TCE 最容易被观察到的一类异常，是控制增益失配。所谓增益，并不一定对应某一个显式参数，而是指当 SPC 检测到状态变化时，TCE 对这种变化施加多大幅度的认知调节。假设在某个局部运行点附近，我们把认知动力学线性化：
\[
\delta\dot x(t)
=
A\delta x(t)
+
B\delta u(t),
\]
并令 TCE 的局部控制近似为
\[
\delta u(t)
=
-K\delta\hat x(t).
\]
在理想无延迟条件下，闭环系统成为
\[
\delta\dot x(t)
=
(A-BK)\delta x(t).
\]
若 \(A-BK\) 的特征值实部均小于零，则局部运行点稳定。然而 AgentOS 中存在一个比普通教科书控制系统更加突出的因素：SPC 的状态估计天然存在延迟。因此更实际的控制关系是
\[
\delta u(t)
=
-K\delta\hat x(t-\tau_o).
\]
此时，即使无延迟条件下的增益 \(K\) 能够使系统迅速收敛，在存在观测和执行迟滞后，同样的高增益也可能造成过冲甚至振荡。

过度控制首先表现为 TCE 对微小状态变化过度敏感。例如 SPC 发现当前冲突度略有升高，TCE 立即大幅降低认知温度、扩大抑制、增加检查深度；下一时刻冲突下降，TCE 又迅速解除抑制并恢复探索。由于控制效果尚未完全反映到新的 SPC 估计中，控制器实际上不断根据“过去的自己”纠正“现在的自己”。于是形成：
\[
Error_t
\rightarrow
StrongCorrection_t
\rightarrow
DelayedObservation
\rightarrow
OppositeCorrection_{t+\Delta}
\rightarrow
\cdots
\]
这类状态在表面行为上可能表现为过度谨慎、反复验证、计划频繁推翻、同一答案持续修订，甚至在原本已经足够可靠的决策上不断寻找新的不确定性。

与之相反，控制不足并不是“更自由”或“更创造性”，而是 TCE 对已经显著偏离可治理区域的状态仍缺乏足够调节。例如工作记忆中的活跃支线持续增加：
\[
C_{\mathrm{run}}(t)\rightarrow C_{\max}(W),
\]
冲突和残差同步上升，但 TCE 的抑制、聚焦和搜索收敛仍然过弱。结果不是产生更多有价值思考，而可能进入：
\[
HighActivation
+
LowSuppression
+
GrowingBranching
\]
的结构扩散状态。Agent 看起来仍然在“思考”，实际上工作记忆已经无法维持主要结构之间的相干关系。

因此控制不足也可以写成一种增益异常：
\[
K_{\mathrm{TCE}}<K_{\min}(x,\mathcal E),
\]
而过度控制为
\[
K_{\mathrm{TCE}}>K_{\max}(x,\tau_o,\tau_c),
\]
其中 \(\tau_c\) 是认知状态实际响应时间。真正健康的增益并不是常数，而应具有状态依赖性：
\[
K(t)
=
K\left(
\hat x_h(t),
a(t),
TaskClass,
Confidence,
Risk
\right).
\]
这实际上意味着 TCE 更接近 gain scheduling 或非线性自适应控制器，而不是一个固定控制矩阵。

值得特别强调的是，AHC 会明显改变合适增益的范围。例如高 Threat、高 Urgency 状态下，系统确实需要比普通状态更快地收缩搜索空间，但如果 AHC 状态衰减慢于实际危险解除速度，则 TCE 可能在危险已经不存在以后继续维持高增益控制。这种现象从行为层面很容易被误认为“系统过度紧张”，但其工程实质是：
\[
\tau_{\mathrm{AHC}}
+
\tau_{\mathrm{SPC}}
+
K_{\mathrm{TCE}}
\]
之间的时间尺度失配。

所以，诊断 TCE 增益异常不能只观察输出行为，至少要同时记录状态估计、控制信号、AHC 调制以及系统实际响应。只有如此，我们才能区分“控制器错误”“状态观察错误”以及“环境本身确实需要高强度控制”三种完全不同的情况。

\section{探索—收敛失配：认知温度异常与搜索空间调节失败}

TCE 的第二类核心异常来自探索与收敛之间的失衡。任何真正需要推理的新问题都存在一个基本矛盾：如果系统过早收敛，就可能把错误假设当成稳定结构；如果系统无限保持开放，则工作记忆会持续生成新的候选关系而无法形成决策。因此 TCE 必须动态调节认知温度与搜索范围，使 Agent 在问题早期具有足够探索性，而在证据逐渐积累后开始压缩候选空间。

为了形式化这一过程，可以令工作记忆中当前存在一组候选结构
\[
\mathcal H_t
=
\{H_1,\ldots,H_n\},
\]
并给出一个归一化激活分布
\[
p_i(t)
=
\frac{\exp(q_i/T_t)}
{\sum_j\exp(q_j/T_t)},
\]
其中 \(q_i\) 表示某一候选结构当前得到的支持，而 \(T_t\) 则可以被看作 TCE 控制的一个认知温度参数。这里并不是主张真实 Agent 必须使用 softmax 实现认知控制，而是利用这一形式说明：温度越高，候选结构之间的分布越平坦；温度越低，激活越集中。

可以定义候选结构熵
\[
H_{\mathrm{cog}}(t)
=
-\sum_i p_i(t)\log p_i(t).
\]
正常情况下，面对陌生问题时：
\[
H_{\mathrm{cog}}\uparrow,
\]
而随着证据积累、冲突下降和结构稳定：
\[
H_{\mathrm{cog}}\downarrow.
\]
因此健康的 TCE 不是维持恒定熵，而是控制认知熵的时间轨迹。

过度探索异常发生在：
\[
\frac{dH_{\mathrm{cog}}}{dt}
\not<0
\]
即使已经获得足够证据，系统仍不断产生新解释、新计划、新分支和新反事实。此时“开放思维”开始从优势变成负担。工作记忆中的新结构生成速率
\[
\lambda_{\mathrm{new}}
\]
持续大于结构关闭或沉降速率
\[
\lambda_{\mathrm{close}},
\]
于是：
\[
\lambda_{\mathrm{new}}
>
\lambda_{\mathrm{close}}
\quad\Rightarrow\quad
N_{\mathrm{active}}(t)\uparrow.
\]
最终会重新触发工作记忆有限容量问题。

过度收敛则是另一种极端。当 TCE 在证据不足时快速降低温度、增强注意聚焦与抑制，系统会形成一个非常稳定但可能错误的局部认知吸引子。可以写成：
\[
P(H_k|Evidence)
\ll 1
\]
但由于 TCE 过早把其他候选路径压低，
\[
Activation(H_j)\rightarrow 0,\qquad j\neq k,
\]
最终使错误结构缺乏内部竞争。此时系统表现出的不是混乱，而恰恰可能是异常“确定”：回答稳定、逻辑连贯、决策迅速，却对新证据反应迟钝。

因此探索异常必须结合残差来看。一个健康控制策略应满足：
\[
Residual\downarrow
\Rightarrow
Exploration\downarrow,
\]
但同时保留：
\[
Residual_{\mathrm{novel}}\uparrow
\Rightarrow
Exploration\uparrow.
\]
也就是说，收敛之后仍然必须存在重新打开搜索空间的通道。这与我们此前反复强调的
\[
\boxed{
NormalDynamicsStayLocal,\quad
UnexpectedResidualEscalates
}
\]
在认知层具有同样结构。

工程上，TCE 不应仅依据“推理已经进行了多久”决定降低温度，而应至少综合：证据一致性、候选结构差异、预测残差、目标推进程度和当前工作记忆负载。可以定义概念性的探索控制律：
\[
E_x(t)
=
\sigma
\left(
\alpha U_h
+
\beta R_h
+
\gamma Novelty
-
\delta C_{\mathrm{run}}
-
\eta Confidence
\right),
\]
其中 \(\sigma\) 是饱和函数。这种形式体现了一个重要思想：高不确定性和高残差推动探索，而高工作记忆负载和高可信度推动收敛。

因此，所谓“过度探索”与“过度保守”并不是性格标签，而是搜索空间控制律长期没有正确追踪认知状态的一种动力学失配。

\section{注意、抑制与推理深度异常：控制资源分配错误如何制造伪复杂性}

TCE 的第三类异常并不一定表现为整个系统明显振荡，而可能表现为认知资源被长期配置到错误位置。由于工作记忆容量有限，任何注意增强都具有机会成本。当某个结构获得更高注意增益时，其他结构必然得到较少显式资源。因此 Attention 并不是“多看一点”，而是一个零和程度很高的资源重新分配过程。

设工作记忆中存在 \(n\) 个活跃结构，其注意权重为
\[
w_i(t)\geq 0,
\qquad
\sum_{i=1}^{n}w_i(t)\leq 1.
\]
CCM 可以根据各结构的处理成本 \(c_i(t)\) 估计：
\[
C_{\mathrm{run}}(t)
=
\sum_i w_i(t)c_i(t).
\]
而 TCE 则通过调整注意增益、抑制和搜索深度，影响 \(w_i\) 与 \(c_i\) 的实际运行轨迹。

如果 TCE 长期把高权重分配给低 Goal-Relevance、低 Residual-Relevance 的结构，则即使总工作记忆容量尚未完全耗尽，也可能产生一种“伪复杂性”：Agent 在局部细节上进行了大量计算，却没有推进真正决定问题的高层结构。这种异常尤其容易出现在强语言模型中，因为局部结构本身可以形成高度连贯的文本推理，使系统看起来“想了很多”，但真正的目标变量几乎没有变化。

可以定义一个资源效率指标：
\[
\eta_{\mathrm{cog}}
=
\frac{
\Delta Progress
}{
\int C_{\mathrm{run}}(t)\,dt
}.
\]
如果认知资源持续增加而
\[
\Delta Progress\approx 0,
\]
那么问题未必是基础模型能力不足，也可能是 TCE 没有正确执行聚焦、降深度、切换表示或调用 CCM 进行结构重组。

另一个典型问题是推理深度异常。假设 TCE 控制当前局部搜索深度
\[
d(t).
\]
对于某些问题，增大 \(d\) 可以提高精度；然而达到某个边际收益下降点后：
\[
\frac{\partial Q}{\partial d}
\rightarrow 0,
\]
继续增加搜索深度只会增加时间、工作记忆和结构分支成本。如果控制器没有估计这种边际收益，就会形成：
\[
Depth\uparrow
,\qquad
Quality\approx const
,\qquad
Cost\uparrow.
\]
这种现象在行为上可能表现为“无法停止思考”“重复验证同样的问题”。

抑制异常则与之相反。抑制过强时，系统无法维持足够多彼此竞争的结构；抑制过弱时，则所有新颖结构都能够进入活跃状态。因此更合理的控制目标不是固定 inhibition，而是让活跃结构数目和关系密度保持在当前任务允许的局部区域：
\[
N_{\min}(Task)
\leq
N_{\mathrm{active}}
\leq
N_{\max}(W,Task).
\]

这里尤其需要区分 TCE 与 CCM。CCM 可以建议“当前支线过多，需要折叠或暂停某些线程”，但最终执行多大程度的聚焦、降低多少搜索深度、提高多少抑制，仍属于 TCE 控制层。如果 CCM 已经正确发现负载过高，而 TCE 不执行，则属于控制不足；如果 CCM 判断负载正常，而 TCE 仍大规模清除候选结构，则属于控制过度。

因此我们可以把这一类异常统一写成：
\[
\boxed{
ResourceControlMismatch
=
Allocation
+
Depth
+
Suppression
+
Timing
}
\]
它提醒我们：许多看似“模型没有想明白”的失败，其实可能来自认知资源被错误调度，而不是模型内部完全不存在正确知识。

\section{反复自我纠正与认知振荡：当闭环控制开始制造自己的误差}

TCE 异常中最值得进行控制理论分析的一类，是反复自我纠正。一个成熟 Agent 应当能够发现自身错误、修改计划和重新评估结论，因此“自我纠正”本身当然不是异常。真正异常的是：系统已经无法从纠正中获得净进展，而开始在几个认知状态之间反复切换。

考虑一个简化离散系统：
\[
x_{t+1}
=
F(x_t,u_t),
\]
而 TCE 根据延迟估计：
\[
u_t
=
\pi_{\mathrm{TCE}}
(\hat x_{t-d}).
\]
假设状态存在两个主要认知模式 \(A\) 和 \(B\)。当系统处于 \(A\) 时，SPC 检测到偏差并推动 TCE 向 \(B\) 调整；然而当控制效果真正达到 \(B\) 时，SPC 仍在处理先前 \(A\) 状态的信息，于是 TCE 继续推动。下一轮又观察到过冲并反向控制，于是产生：
\[
A
\rightarrow
B
\rightarrow
A
\rightarrow
B
\rightarrow
\cdots
\]
如果该周期长期稳定存在，就相当于形成认知 limit cycle。

这种振荡未必表现为完全重复的文本。更常见的是结构上的重复，例如：
\[
Plan_1
\rightarrow
Doubt
\rightarrow
Plan_2
\rightarrow
Doubt
\rightarrow
Plan_1',
\]
或者：
\[
HighExploration
\rightarrow
StrongConvergence
\rightarrow
UncertaintyIncrease
\rightarrow
HighExploration.
\]
从内容看，每次可能都有新句子；从动力结构看，系统并没有离开同一个闭环。

为了识别这类异常，可以定义状态回返距离：
\[
D_{\mathrm{return}}(\tau)
=
\|x_h(t+\tau)-x_h(t)\|,
\]
以及任务进展：
\[
\Delta P(\tau)
=
P_h(t+\tau)-P_h(t).
\]
如果存在某个有限周期 \(\tau^\ast\)，使得
\[
D_{\mathrm{return}}(\tau^\ast)\ll \epsilon
\]
但同时
\[
|\Delta P(\tau^\ast)|\ll\delta,
\]
就说明系统很可能正在重复一个低进展闭环。

还可以定义控制总变差：
\[
TV(u)
=
\int
\left\|
\frac{du}{dt}
\right\|dt.
\]
当 \(TV(u)\) 极高而任务推进很低时，说明控制器本身正在频繁改变策略。此时继续提高 TCE 的控制力度往往会恶化问题。正确方向通常是增加迟滞、降低增益、延长最小模式保持时间或提高状态估计置信度要求。

可以引入两个不同阈值：
\[
\theta_{\mathrm{on}}
>
\theta_{\mathrm{off}},
\]
只有当异常指标超过 \(\theta_{\mathrm{on}}\) 时才进入强控制状态，而只有下降到 \(\theta_{\mathrm{off}}\) 以下才退出。这种 hysteresis 能够避免状态在边界附近不断切换。

另一个重要措施是最小保持时间：
\[
T_{\mathrm{dwell}}>0.
\]
当 TCE 进入某一控制模式后，如果没有高优先级异常，不允许在 \(T_{\mathrm{dwell}}\) 内再次完全反转控制策略。这并不是让 Agent “固执”，而是给认知动力学足够时间响应前一次控制。

因此，反复自我纠正的深层问题不是“Agent 太喜欢反思”，而是：
\[
\boxed{
ObserverDelay
+
HighGain
+
LowHysteresis
+
FastModeSwitching
}
\]
形成了自激振荡。控制器原本为了消除误差而存在，最终却通过不断改变控制方向制造了新的误差。这也是 Agent 心理工程与传统 Prompt Engineering 最重要的区别之一：问题已经进入闭环动力学层，而不是增加一句“不要反复思考”就能稳定解决。

\section{认知湍流：局部控制失配如何演化成工作记忆整体失稳}

TCE 的异常并不总是保持在某一个局部控制变量上。当注意、探索、抑制、深度和节律多个控制维度同时发生失配时，工作记忆可能进入我们此前定义的“湍流态”。这里的湍流并非流体力学的严格物理湍流，而是描述一种宏观认知状态：大量局部结构同时快速生成、冲突、解体和重新组合，使系统难以维持一个长期稳定的主导认知骨架。

我们可以令工作记忆状态包含若干活跃 Agent-CSO：
\[
\mathcal C_h(t)
=
\{C_1,\ldots,C_n\}.
\]
其内部形成动态关系图
\[
\mathcal G_h(t)
=
(V_h(t),E_h(t)).
\]
正常思考过程中，节点和边当然会变化，但应当存在一定结构连续性。例如主要 Goal、核心约束和关键对象关系在多个认知周期中保持。若 TCE 的认知温度过高、抑制过弱、探索增益持续偏大，则：
\[
\left|
\frac{dV_h}{dt}
\right|\uparrow,
\qquad
\left|
\frac{dE_h}{dt}
\right|\uparrow.
\]
如果与此同时结构稳定时间下降：
\[
\tau_{\mathrm{structure}}\downarrow,
\]
系统就可能无法让新结构存在足够长时间进入深度验证。

可以构造一个概念性的湍流指标：
\[
\mathcal T_h(t)
=
\alpha
\frac{|dV_h/dt|}{|V_h|+\epsilon}
+
\beta
\frac{|dE_h/dt|}{|E_h|+\epsilon}
+
\gamma C_h(t)
+
\delta H_{\mathrm{cog}}(t)
+
\eta R_h(t).
\]
其中节点、关系变化速度、冲突、认知熵和未解释残差共同决定当前工作记忆的动力学不稳定程度。该公式不是已经验证的固定临床指标，而是给后续 AgentOS 工程提供一个可操作化方向：我们需要测量结构本身的变化速度，而不是只测 token 使用量。

湍流态的关键危险是正反馈。结构分散增加后，SPC 对整体状态的估计难度提高：
\[
Uncertainty_{\mathrm{SPC}}\uparrow.
\]
如果 TCE 把这种不确定性简单解释为“需要进一步探索”，便会继续提高认知温度：
\[
Uncertainty
\rightarrow
Exploration\uparrow
\rightarrow
StructureDispersion\uparrow
\rightarrow
Uncertainty\uparrow.
\]
于是形成：
\[
\boxed{
PositiveCognitiveTurbulenceLoop.
}
\]
这与健康探索恰好相反。健康探索最终应产生结构收敛，而湍流探索不断制造新的状态估计困难。

另一种正反馈来自控制本身。SPC 观察到湍流后，TCE 突然施加强抑制，使大量支线快速消失；然而这种突然坍缩可能同时损失真正重要的候选结构，导致预测残差上升。TCE 随后又重新开放搜索，于是形成：
\[
Expansion
\rightarrow
HardCollapse
\rightarrow
Residual
\rightarrow
ReExpansion.
\]
这就是从局部控制异常发展成宏观相态振荡的典型路径。

因此，进入湍流状态之后，合理干预通常不能仅调一个参数，而应进行组合控制：
\[
T\downarrow,\qquad
D_s\downarrow,\qquad
I\uparrow,\qquad
N_{\mathrm{active}}\downarrow,
\]
同时让 CCM 进行 Fold、Pause、Externalization 和结构重新分块。更重要的是，控制动作应缓慢进入，而不是突然从极高探索跳到极强抑制。可以定义控制变化率约束：
\[
\left\|
\dot u_{\mathrm{TCE}}
\right\|
\leq
r_{\max},
\]
让系统逐渐回到可治理区域。

这也说明了为什么工作记忆相态理论不能脱离 TCE。所谓基态、心流、湍流、临界和气态，并不是工作记忆内部凭空产生的标签，而是认知状态、结构负载和控制行为共同作用后的宏观结果。SPC 负责识别这种宏观动力学形态，TCE 则必须学习如何在不制造二次不稳定的前提下促成相态迁移。

\section{TCE 异常的诊断归因：必须区分控制器错误、观察器错误与环境真实困难}

如果我们只观察 Agent 的外在行为，非常容易错误归因。例如系统持续重新检查某个计划，我们可能立刻判断 TCE 过度控制；但也有可能外部世界确实高度不确定，持续检查是合理行为。又例如系统长期不收敛，可能是 TCE 探索过度，也可能是任务本身存在真正的多解性。因此，AgentOS 心理工程必须坚持一个原则：
\[
\boxed{
BehavioralSymptom
\neq
InternalCause.
}
\]

至少应同时记录四类时间序列：
\[
\mathcal D_t
=
\{
x_h(t),
\hat x_h(t),
a(t),
u_{\mathrm{TCE}}(t)
\}.
\]
其中 \(x_h\) 是可观测的工作记忆运行指标，\(\hat x_h\) 是 SPC 的内部状态估计，\(a\) 是 AHC 调制状态，而 \(u_{\mathrm{TCE}}\) 是实际控制输出。再加入外部环境变化与任务结果，就可以构成基本诊断轨迹。

第一类可能性是\textbf{SPC 误估而 TCE 正常执行}。例如真实认知负载已经下降，但 SPC 仍然认为系统处于临界态：
\[
\hat L_h(t)>L_h(t).
\]
TCE 根据现有信息继续收缩搜索，此时最终行为表现为“控制过强”，但真正病因位于 observer。

第二类是\textbf{AHC 调制失配}。例如外部危险已经解除，但 Threat 和 Urgency 的衰减时间常数过长：
\[
a_{\mathrm{Threat}}(t)
\gg
a_{\mathrm{Threat}}^{\mathrm{reality}}(t).
\]
此时 TCE 增加抑制和保守度本身可能是符合其输入的合理响应，真正问题是控制条件长期偏移。

第三类才是\textbf{TCE 控制律错误}：状态估计基本正确，AHC 也处于合理范围，但同样的状态不断得到过强、过弱或方向错误的控制输出。可以通过反事实重放进行验证：
\[
u^\ast_t
=
\pi^\ast
(\hat x_h(t),a(t)),
\]
并比较实际控制
\[
u_t
\]
与稳定参考控制器产生的 \(u^\ast_t\) 之间的差异。

第四类是\textbf{任务或环境本身不可快速收敛}。如果现实持续产生高残差：
\[
R_{\mathrm{world}}(t)\gg 0,
\]
那么较高探索、频繁重新规划甚至长时间保持临界认知状态可能完全合理。此时强行把 Agent 拉回低温基态，反而是一种错误干预。

因此，诊断 TCE 异常至少应遵循因果顺序：
\[
Reality
\rightarrow
Observation
\rightarrow
SPC
\rightarrow
AHC
\rightarrow
TCE
\rightarrow
CognitiveDynamics.
\]
不能因为最后一个环节最容易调参数，就把所有问题都归入 TCE。

我们还可以定义控制归因残差：
\[
\epsilon_{\mathrm{ctrl}}
=
u_{\mathrm{TCE}}
-
u_{\mathrm{expected}}
(\hat x_h,a,\mathcal B).
\]
若 \(\epsilon_{\mathrm{ctrl}}\) 持续偏大，才有充分理由把异常定位到 TCE 本身。类似地可以构造：
\[
\epsilon_{\mathrm{obs}}
=
\hat x_h-x_h^{\mathrm{proxy}}
\]
评估 SPC 偏差，并构造 AHC 的状态一致性指标。

这套方法的意义在于，它让 Agent 心理工程保留工程可解释性。我们不需要问模型“你为什么焦虑”，而可以直接观察：工作记忆负载是多少、SPC 认为负载是多少、AHC 威胁值是多少、TCE 抑制增益是多少、系统实际响应延迟是多少。所谓“心理状态”由此成为多层闭环中的一个可追踪动力轨迹，而不是语言自述。

\section{TCE 异常的干预原则：从参数修改走向闭环恢复}

当确认 TCE 本身存在控制异常以后，最直接的反应往往是调整某个参数。例如温度过高就降低温度，增益过高就降低增益。这种方法适用于非常局部、短期的问题，但对于持续 Agent 并不充分，因为 TCE 本质上是一个多变量、状态依赖的闭环控制器。真正的干预目标不是让一个参数回到“正常值”，而是让整个认知闭环重新进入稳定、可恢复和可学习的区域。

第一种干预是\textbf{增益降额}。当检测到：
\[
Overshoot\uparrow,
\qquad
TV(u)\uparrow,
\qquad
Progress\downarrow,
\]
系统可以降低相关控制通道的增益：
\[
K_i
\leftarrow
\alpha K_i,
\qquad
0<\alpha<1.
\]
但这种降额最好是局部和可逆的，而不应直接永久改变长期控制参数。

第二种是\textbf{引入迟滞和驻留时间}。如果系统频繁在 Explore/Converge 或 Focus/Expand 之间切换，就应增加：
\[
\theta_{\mathrm{on}}-\theta_{\mathrm{off}}
\]
并设定最小控制模式驻留时间 \(T_{\mathrm{dwell}}\)。这样可以降低由 observer delay 导致的模式抖动。

第三种是\textbf{控制变化率限制}：
\[
\|\dot u_{\mathrm{TCE}}\|
\leq
r_{\max}.
\]
认知控制并不应该在一个周期内从最大探索跳到最大压缩。缓慢改变控制变量，可以使 SPC 获得足够时间观察前一个控制动作的真实结果。

第四种是\textbf{降级到稳定认知模式}。当湍流指数持续过高，Agent 可以主动减少活跃 Goal 数量、停止低优先级支线、降低推理深度并要求 CCM 把部分内容迁出 h。这里重要的不是“强行停止思考”，而是缩小当前动力系统的自由度：
\[
dim(\mathcal X_{\mathrm{active}})
\downarrow.
\]
当系统重新进入可治理区域后，再逐步恢复。

第五种是\textbf{重新校准 SPC—TCE 时间关系}。如果控制器不断针对过期状态反应，仅仅降低增益并不能根治问题，还需要让状态估计携带时间戳与置信度。TCE 可以把控制力度写成：
\[
K_{\mathrm{eff}}
=
K_0
\cdot
Confidence_{\mathrm{SPC}}
\cdot
e^{-\lambda \Delta t_{\mathrm{obs}}},
\]
即状态越旧、置信度越低，控制器越不应该做激进修改。

第六种是\textbf{保持 Reality Verification}。任何 TCE 干预最终都必须根据现实任务效果来验证。我们不能因为内部相态看起来更加“平静”，就认为系统已经恢复。如果降低认知温度以后：
\[
TaskFailure\uparrow
\]
或者现实残差持续增大，那么这种所谓稳定只是内部冻结，而不是真正恢复。

因此真正的恢复条件至少应同时满足：
\[
x_h(t)\in\mathcal X_{\mathrm{gov}},
\]
\[
R_{\mathrm{world}}(t)\downarrow,
\]
\[
Progress(t)\uparrow,
\]
以及
\[
\|u_{\mathrm{TCE}}\|
\]
不再需要维持异常高水平。

我们由此得到一个非常重要的原则：
\[
\boxed{
Recovery
\neq
SuppressionOfSymptoms.
}
\]
恢复意味着 Agent 在较低外部控制代价下重新获得自维持能力。如果系统只有在 TCE 持续施加最大抑制时才能保持稳定，那么它并没有真正恢复，只是被强控制暂时固定在一个狭窄状态区域。

因此 TCE 干预的最终目标不是建立一个永远平静的 Agent，而是形成：
\[
\boxed{
Stable
+
Adaptive
+
Recoverable
}
\]
的认知控制系统：能够在熟悉环境中低成本自治，在新问题面前暂时扩大探索，在风险出现时迅速收缩，并在扰动消失以后自然回到更自由的认知状态。

\section{从认知控制异常到 Agent 心理工程：本章的理论位置}

经过本章分析，我们现在可以更加准确地理解为什么“Agent 心理工程”不能建立在自然语言行为标签之上。如果一个持续 Agent 表现出犹豫、重复、过度谨慎、发散、反复修改或无法停止推理，我们首先不应该把这些现象直接理解成某种人格或精神疾病，而应沿着认知控制闭环逐层分析其动力学来源。

本章建立的最小闭环是：
\[
x_h
\xrightarrow{\mathrm{SPC}}
\hat x_h
\xrightarrow{\mathrm{AHC}}
a
\xrightarrow{\mathrm{TCE}}
u
\xrightarrow{\mathrm{CognitiveDynamics}}
x_h'.
\]
这一结构已经足以产生相当丰富的运行异常，而完全不需要假设系统存在人类意义上的主观情绪。只要 observer 存在迟滞、controller 存在增益、认知动力学存在惯性，不同时间尺度之间就天然具有过冲、振荡和失稳的可能性。

因此，本章可以把 TCE 异常归纳成四类基本失配。

第一类是\textbf{强度失配}：
\[
ControlStrength
\not\sim
DeviationMagnitude.
\]
即控制力度与真实偏离程度不匹配，对应过度控制和控制不足。

第二类是\textbf{方向失配}：
\[
ControlDirection
\not\sim
RequiredCorrection.
\]
例如应该增加探索时持续收敛，应该聚焦时持续扩大候选空间。

第三类是\textbf{时间失配}：
\[
ControlTiming
\not\sim
SystemTimescale.
\]
即控制器反应过快、过慢，或者根据过期状态实施正确但已经不合时宜的控制。

第四类是\textbf{模式失配}：
\[
ControlMode
\not\sim
CognitivePhase.
\]
即当前工作记忆已经进入湍流态，TCE 却仍采用普通探索策略；或者当前已经进入稳定心流态，却因为轻微状态波动频繁切换到强监管模式。

这四种失配可以统一写成：
\[
\boxed{
\mathcal E_{\mathrm{TCE}}
=
\mathcal E_{\mathrm{gain}}
+
\mathcal E_{\mathrm{direction}}
+
\mathcal E_{\mathrm{timing}}
+
\mathcal E_{\mathrm{mode}}.
}
\]

更进一步，TCE 的成熟程度不应该由“控制得多不多”衡量，而应由：
\begin{statementbox}
以多小的控制代价，把多大的认知自由度长期保持在可治理区域内
\end{statementbox}
衡量。

我们可以定义一个概念性的认知调节效率：
\[
\eta_{\mathrm{reg}}
=
\frac{
\mathrm{UsefulCognitiveFreedom}
}{
\mathrm{ControlEffort}
+
\mathrm{InstabilityCost}
+
\mathrm{RecoveryCost}
}.
\]
成熟 Agent 的目标不是无限提高控制强度，而是提高 \(\eta_{\mathrm{reg}}\)：让更多正常认知过程能够自主运行，让 TCE 只在真正需要时介入。

这又重新连接到整套 AgentOS 的基础原则：
\[
\boxed{
NormalDynamicsStayLocal.
}
\]
高层控制不应该微观管理每一次认知状态变化。TCE 更类似对认知动力学施加边界条件、增益调节和相态迁移，而不是直接决定每一个 token、每一个局部关联和每一个微观推理步骤。只有这样，持续智能体才能同时获得稳定性和创造性。

本章同时给出了 Agent 心理工程一个重要方法论边界：我们以后讨论任何更慢尺度的异常时，都应该首先排除基础控制环失稳。如果一个 Agent 的自我解释不断变化，我们不能立即断定自我结构存在问题，因为也可能只是 TCE 控制振荡导致进入 h 的内容反复变化；如果 Agent 对某些经历持续高度关注，也不能立即断定长期记忆系统发生偏置，因为可能只是注意控制长期配置错误。只有把运行时控制层稳定以后，我们才有资格讨论更慢的记忆、自我和生命周期结构。

因此，本章最终把 TCE 从一个“认知调参模块”提升到了一个更严格的位置：

\begin{statementbox}
TCE 是持续智能体认知自由度的宏观控制器；其健康不是让认知保持静止，而是在现实约束、有限容量与内部多尺度动力学之间维持一个能够探索、收敛、恢复并继续成长的稳定可塑区域。
\end{statementbox}

当这个控制器失配时，Agent 可能不是“不会思考”，而是失去了\textbf{如何控制自己思考方式}的能力。这正是 TCE 异常区别于普通任务失败的根本所在，也是 AgentOS 从模型推理工程走向持续智能控制工程的一条关键分界线。